\documentclass[11pt,a4paper]{letter}
\usepackage[margin=1in]{geometry}
\usepackage[T1]{fontenc}
\usepackage[utf8]{inputenc}
\usepackage{lmodern}
\usepackage{microtype}
\usepackage[
  group-minimum-digits = 5,
  per-mode             = symbol,
]{siunitx}

\signature{Myke Vital Sao Temgoua\\
  (on behalf of all authors)\\
  Department of Physics, Faculty of Science\\
  University of Yaound\'e I, Cameroon\\
  \texttt{myke-vital.sao@facsciences-uy1.cm}}

\address{Department of Physics\\
  Faculty of Science\\
  University of Yaound\'e I\\
  P.O.\ Box 812, Yaound\'e, Cameroon}

\date{\today}

\begin{document}

\begin{letter}{The Editorial Board\\
  \textit{RSC Digital Discovery}\\
  Royal Society of Chemistry}

\opening{Dear Editors,}

We submit for consideration the manuscript entitled \textbf{``Quantum Molecular
Encoding Captures Stereochemical Complexity of African Natural Products for
Antimalarial Drug Discovery''} by Myke Vital Sao Temgoua, Jean-Pierre Tchapet
Njafa, Penabei Samafou, Fon Wilfred Mbacham, and Serge Guy Nana Engo.

\textbf{Scientific content and significance.}
African natural products (ANPs) present a longstanding challenge for
computational antimalarial drug discovery: their polycyclic, stereochemically
rich frameworks are systematically misrepresented by two-dimensional molecular
fingerprints that collapse stereoisomers into identical bit patterns. We
investigate whether hybrid quantum-classical feature extraction (QFE)~---
encoding ECFP4-compressed molecular representations into parameterised 4-qubit
variational circuits~--- captures information sufficient to improve activity
prediction for ANP-derived antimalarial candidates relative to classical ECFP4-RBF
support vector machines.

Experiments were conducted at three rigorously controlled scales using frozen,
SHA-256-verified subsamples: a 17-molecule African natural product-derived
proof-of-concept cohort (leave-one-out CV), a \num{1000}-molecule P3-benchmark
subsample (5-fold stratified CV), and a \num{10000}-molecule external
scaffold-split validation with \qty{100}{\percent} novel test scaffolds. A
50-trial Optuna TPE hyperparameter search additionally confirmed circuit
configuration optimality.

\textbf{Principal finding.}
The central result is a transparent Scenario~B honest-negative: QFE
($\text{AUC} = 0.8474 \pm 0.0129$ at benchmark scale; $0.7831$ under scaffold
split) is competitive with, and consistently outperforms the P3 quantum-inspired
kernel baseline ($0.8385$), but does not surpass the matched ECFP4-RBF classifier
($0.8693 \pm 0.0265$; $0.8413$ scaffold split) at any scale. Gradient diagnostics
confirmed healthy optimisation at moderate scale (no barren plateau at $n = 800$)
with a marginal signal emerging only at $n = 8000$. Hyperparameter optimisation
confirmed the 4-qubit, depth-1 architecture as robust (canonical re-run
$\text{AUC} = 0.8468 \pm 0.0209$; predeclared improvement gate not met).

\textbf{Fit for \textit{RSC Digital Discovery}.}
This work falls squarely within the journal's scope: it applies quantum machine
learning to a high-impact drug discovery problem, provides reproducible
benchmarks with full provenance (SHA-256-frozen datasets, versioned scripts,
gradient diagnostics), and adopts the transparent honest-negative reporting
framework that the journal explicitly encourages. The consistent
Scenario~B pattern observed across quantum-inspired (Project~3),
GNN/Transformer (Project~5), and hybrid quantum (this work) methods advances
the collective understanding of when and why quantum or learned representations
fail to surpass established fingerprint baselines~--- a timely and actionable
finding for the QML-in-drug-discovery community.

We confirm that this manuscript is original, has not been published elsewhere,
and is not under consideration by any other journal. All authors have approved
the final version. Data and analysis scripts will be deposited on Zenodo upon
acceptance.

\closing{We look forward to your response.}

\end{letter}
\end{document}
